\section{Changes with respect to the published version} \label{app:changes}

\begingroup\raggedright
The base text of this revision is arXiv:1206.2251v2. The version published in Duke Math.\ J.\ \textbf{163} (2014) is not in the repository \cite{Lean}, and differences between it and arXiv:1206.2251v2, if any, are not recorded here.

The numbering of sections, statements and displays is that of arXiv:1206.2251v2. The new statements (Remarks 1.5, 3.9 and 4.1, Lemma 6.6) come last in their sections, and the new displays are unnumbered or tagged, (2.4$'$) and (3.20$'$), so no number changes. Each line below gives the location in this revision and the change in short (old $\to$ new). Obvious spelling slips in the English text (such as ``statics'', ``stretagy'' and ``Isotrpic'') were corrected without being listed.

\medskip
\noindent\textbf{Formalization.} Every theorem, proposition, lemma and corollary of this revised paper that the proof of Theorem \ref{main result} uses is formalized in \cite{Lean} in the form stated here, and so are the results of \cite{EKYY1}, \cite{EKYY2}, \cite[Lemma 4.2]{EYY2} and \cite{EYY} in the forms stated here (Lemma \ref{lm:msc}, Lemma \ref{lem: old}, Lemma \ref{lem: ex edge un}, Lemma \ref{lem: edge un str} and the steps of the proof of Theorem \ref{thm: new rigidity} that follow \cite[Section 5]{EYY}). The Lean names, all in the namespace \texttt{WEdgeNS}, are:
\begin{itemize}
\item Theorem \ref{main result}: \texttt{E1}, \texttt{E1\_not\_tight}, \texttt{E1\_no\_limit} (necessity) and \texttt{E2} (sufficiency);
\item Lemma \ref{lm:msc}: \texttt{norm\_msc\_eq\_inv\_and\_lt\_one}, \texttt{msc\_bounds};
\item Lemma \ref{lem: old}: \texttt{localLaw}, \texttt{norm\_le\_two\_add}, \texttt{delocalization}, \texttt{rigidity}, \texttt{outOfSpectrum};
\item Lemma \ref{lem: ex edge un}: \texttt{edge\_universality\_largeQ}; Theorem \ref{thm: new rigidity}: \texttt{rigidity\_smallQ}; Theorem \ref{thm: new edge un}: \texttt{edge\_universality\_smallQ}; Lemma \ref{lem: WBGij}: \texttt{offDiag\_sq\_le\_smallQ};
\item Lemma \ref{lem: match}: \texttt{exists\_momentMatched}; Lemma \ref{lem: four moments}: \texttt{exists\_law\_four\_moments}; Remark \ref{rem: four moments}: \texttt{sq\_moment\_three\_add\_one\_le}; Lemma \ref{lem:p moment}: \texttt{moment\_comparison}; Lemma \ref{lem:green}: \texttt{green\_comparison\_smallQ}; Lemma \ref{lem: edge un str}: \texttt{edge\_universality\_of\_inputs};
\item Lemma 6.4: \texttt{norm\_evalW\_wordP\_le\_of\_good} on the event of \texttt{highProb\_goodEvent}; Lemma \ref{lem:S_expansion}: \texttt{expectation\_expansion}, \texttt{expectation\_expansion\_inv}; the iteration, \eqref{HHtt} and (6.37): \texttt{comparison\_iterate\_sum}, \texttt{comparison\_iterate}; Lemma \ref{lem:words}: \texttt{norm\_integral\_evalF\_iterP\_le}.
\end{itemize}
The following are not formalized, and the proof of Theorem \ref{main result} does not use them: Remark \ref{rem: TW limit} (the Tracy--Widom limit, which rests on \cite{TR1, TR2}); the Hermitian case and the case $k\ge2$ mentioned in Remark \ref{rem: scope}; the example of Remarks \ref{example} and \ref{rem: additive error}; the sentence on the error $N^{-\delta'}$ at the end of Remark \ref{rem: additive error}; Theorem \ref{thm: old edge un} in its stated form, with the bounded support condition of Definition \ref{cond:bdd} (it is used only through Lemma \ref{lem: ex edge un}); Remark \ref{genedge} and \eqref{twa}; the probabilistic form of Lemma \ref{lem: ex edge un} mentioned after it; the form of Theorem \ref{thm: new rigidity} under the bounded support condition of Definition \ref{cond:bdd} mentioned after it; and the heuristic Example 6.1.

\subsection*{Title page, abstract and bibliography}
\begin{itemize}
\item Abstract: ``Tracy-Widom law holds if and only if \eqref{crit}; the same criterion holds for Hermitian Wigner matrices'' $\to$ for real symmetric Wigner matrices, \eqref{crit} is necessary and sufficient for the edge universality of $\lambda_N$; under \eqref{crit} the distribution of $N^{2/3}(\lambda_N-2)$ is that of the GOE up to a shift $N^{-\delta}$ and an additive error $o(1)$ uniform in $s$, hence $F_1$ by \cite{TR1, TR2}; otherwise it is not tight; the proof is formalized in Lean 4 \cite{Lean}. The Hermitian sentence is removed.
\item Revision line after the title: new.
\item Bibliography: \cite{EKYY1} ``arXiv:1103.1919v2 (2011)'' $\to$ ``Ann.\ Probab.\ \textbf{41} (2013); arXiv:1103.1919v5''; \cite{EKYY2} ``arXiv:1103.3869v2 (2011)'' $\to$ ``arXiv:1103.3869v3''; \cite{EYY}: arXiv:1007.4652v7 added; \cite{Lean}: new entry, the Lean formalization, with a breakable URL.
\end{itemize}

\subsection*{Section 1}
\begin{itemize}
\item Paragraph after (1.3), before \eqref{crit}: ``the joint probability density of the $k$ largest eigenvalues \dots\ weakly converge to that of Gaussian ensembles'' $\to$ ``the rescaled largest eigenvalue of a real symmetric Wigner matrix has asymptotically the same distribution as for the GOE, and hence the Tracy--Widom distribution $F_1$''.
\item Definition \ref{def: SWigner}: new last sentence, Theorem \ref{main result} concerns the real symmetric case.
\item Theorem \ref{main result}, \eqref{def: ranvec}--\eqref{nec}, and the new paragraph after it: real symmetric case and $\lambda_N$ only; sufficiency: ``for any fixed $k$ the joint distribution of the $k$ rescaled largest eigenvalues converges to the GUE (GOE) limit'' $\to$ the GOE comparison \eqref{def: ranvec} with a shift $N^{-\delta}$ and an additive error $\epsilon_N\to0$ uniform in $s$; necessity: ``\eqref{crit} does not hold'' is read as ``$s^4\p(|x_{12}|\ge s)$ does not tend to $0$'', and the consequences ``not tight'' and ``no limit distribution'' are stated. The paragraph after the theorem defines $\p^{GOE}$ and $F^{GOE}_N$, fixes the order of the quantifiers and derives the two consequences from \eqref{nec}.
\item Remark \ref{example}: density ``decaying as $|x|^{-5}\log|x|$'' (for which \eqref{crit} fails) $\to$ $c|x|^{-5}(\log|x|)^{-1}$ on $|x|\ge e$, normalised to variance $1$, with $s^4\p(|x_{12}|\ge s)\to0$ and $\E x_{12}^4=\infty$; new sentence: for this law, \eqref{def: ranvec} fails with $\epsilon_N=N^{-\delta}$ (Remark \ref{rem: additive error}).
\item After Remark \ref{example}, the two key observations: ``entries have a bounded support of size $N^{-\epsilon}$'' $\to$ ``entries are almost surely bounded by $N^{-\epsilon}$''; ``first four moments coincide'' $\to$ ``first four moments of the off-diagonal entries and first two moments of the diagonal entries coincide''. The organisation paragraph ends with a pointer to this appendix.
\item Remark \ref{rem: scope}: ``we prove only the real symmetric case with $k=1$; the general case can be proved analogously'' $\to$ the Hermitian case and $k\ge2$ are expected to follow from the same method but are not proved here.
\item Remark \ref{rem: TW limit} (new): deduction of $\p(N^{2/3}(\lambda_N-2)\le s)\to F_1(s)$ from \eqref{def: ranvec} and \cite{TR1, TR2}; not formalized. After it, a new unnumbered paragraph ``Formalization''.
\end{itemize}

\subsection*{Section 2}
\begin{itemize}
\item Definition \ref{cond:bdd}, \eqref{eq:cond_bdd} and new (2.4$'$): ``$|h_{ij}|\le q^{-1}$ with probability larger than $1-e^{-N^c}$ for some $c>0$; usually $N^\phi\le q\le N^{1/2}(\log N)^{-1}$'' $\to$ $\p(|h_{ij}|>q^{-1})\le e^{-N^c}$ with $c$ a fixed parameter, plus the tail-variance condition (2.4$'$) for $i\ne j$; the almost sure form is defined; $q\le\sqrt2N^{1/2}$.
\item Paragraph after (2.4$'$): ``throughout the paper, we will not consider the case it happens'' $\to$ the event $\{|h_{ij}|>q^{-1}\}$ is removed by an explicit truncation in the proof of Lemma \ref{lem: old}, and it does not occur in Sections 5 and 6, where the entries are bounded almost surely.
\item Definition 2.7, \eqref{def:gamma}: new sentence, $\gamma_j$ is the smallest solution, so $\gamma_N=2$.
\item Remark 2.8: new sentences on the dependence of constants on the fixed parameters and on $N$-dependent distributions.
\end{itemize}

\subsection*{Section 3}
\begin{itemize}
\item Lemma \ref{lem: old}, (3.1)--\eqref{3838}: now ``fix $\zeta>0$, $C_0\ge0$, $c>0$'', with \eqref{eq:diag} and Definition \ref{cond:bdd}, $\zeta$-high probability (was $3$-high probability), constants depending on $(\zeta,C_0,c)$ only; \eqref{scm}, \eqref{Gij estimate} and (3.6) with $\bigcap$ (printed $\bigcup$); \eqref{deloc} for every unit eigenvector; parts 4 and 5 for $q\ge N^\phi$, $\phi>1/3$; in part 5 the offset $3$ of \eqref{3737} becomes a constant $K$, $a>12$ becomes $a>12\zeta$, and \eqref{3838} holds simultaneously for all $E$ in \eqref{3737}.
\item Proof of Lemma \ref{lem: old} (after \eqref{3838}), including \eqref{tyzz}: rewritten in five steps: an explicit truncation reduces to entries bounded almost surely; \cite{EKYY1} is run at level $\zeta$ with diagonal variances bounded by $C_0$ (no rescaling by $1+O(N^{-1})$), extended to $2\le\eta\le10$, and made uniform in $z$ by a lattice; delocalization and rigidity (including $j=N$) are derived from parts 1 and 2; the offset of part 5 is a constant $K_0$ (the printed $3$ conflates two constants $C_1$ of \cite{EKYY1}); \eqref{tyzz}: $O(q^{-1})+(\varphi^{a/2-2}N\eta)^{-1}$ $\to$ $C_o(q^{-1}+\varphi^{2\zeta-a/2}(N\eta)^{-1})$; ``$\eta\ge\varphi^{O(1)}N^{-2/3}$'' (wrong direction) $\to$ $N\eta\le\varphi^{K+a-C/2}N^{1/3}$.
\item Remark \ref{prozeta}: ``we only use $1\le\zeta\le3$'' $\to$ Lemma \ref{lem: old} holds for every fixed $\zeta$, and $\zeta\in\{1,3,4\}$ is used, $\zeta=4$ for \eqref{3838T}.
\item Theorem \ref{thm: old edge un}, \eqref{qset}, and the new paragraph after it: cited as \cite[Theorem 2.7 and Proposition 6.4, (6.15) at $f=0$]{EKYY2}; $\delta$ and $N_0$ independent of $s$; the variance convention of \cite{EKYY2} is noted; in the stated form the theorem is cited and not used.
\item Remark \ref{genedge}, \eqref{twa}: new sentence, not used (see \eqref{segg}) and not formalized.
\item Lemma \ref{lem: ex edge un} and the sentence after it: ``the same condition except $\E|h_{ii}|^2\le C_0/N$'' $\to$ \eqref{eq:diag} with $C_0\ge0$ and $|h_{ij}|\le N^{-\phi}$ almost surely, $1/3<\phi\le1/2$, $\delta$ and $N_0$ depending on $(C_0,\phi)$, all $N\ge N_0$ and all $s$; the probabilistic form is not proved here.
\item Theorem \ref{thm: new rigidity}, \eqref{rigiditynew}--\eqref{rigiditynew2}, and the new paragraph after it: bounded support with $q=N^\phi$ $\to$ $|h_{ij}|\le N^{-\phi}$ almost surely; \eqref{rigiditynew} with $\bigcap_j$ (printed $\bigcup_j$); \eqref{rigiditynew2} for all $E\in\mathbb R$ (was $|E|\le5$); the form under Definition \ref{cond:bdd} is not proved here.
\item Theorem \ref{thm: new edge un}, \eqref{qset1}: entries bounded almost surely; $\delta$, $N_0$ depend only on $(C_1,\phi)$; all $N\ge N_0$, all $H^{\V}$ and all $s$.
\item Lemma \ref{lem: WBGij}, \eqref{res: WBGij}: entries bounded almost surely; for every $0<c<1$ there are $C$ and $N_0$ depending on $(C_1,\phi,c)$, for all $i\ne j$.
\item Proof of Lemma \ref{lem: ex edge un}, \eqref{jhw}--\eqref{eq:tel_1}: rewritten. ``A simple extension of \cite[Theorem 2.7]{EKYY2}'' $\to$ a derivation from Lemma \ref{lem: edge un str} with rate $N^{-\kappa+C\epsilon}$, $\kappa=\phi-1/3$, its hypotheses checked from Lemma \ref{lem: old}; the Claim with explicit constants, $F\in C^4$ (constant renamed $C_F$) and $E_1,E_2$ in either order; reduction to the off-diagonal parts by the Ward identity, new display (3.20$'$); the local law for every $H_\gamma$ with $q'=N^{(1/3+\phi)/2}$ (the interpolating matrices do not satisfy the support condition with $q$); \eqref{eq:S_bound}--\eqref{8800} with explicit constants and the entry bound $N^{\epsilon}q^{-1}$; the case $a\ne b$, \eqref{eq:S-T 1}, follows \cite[(6.41)--(6.97)]{EKYY2} at $f=0$ with the changes listed there; the case $a=b$, \eqref{eq:short_exp}--\eqref{eq:S-T 2}, uses that the linear term has mean zero, the bad-event bound \eqref{eq:diag_bad} (new label) and only an upper bound on $\E v^2$ (the diagonal second moments need not match); the integrated functional and the diagonal parts are treated explicitly; telescoping \eqref{eq:tel_1}.
\item Remark \ref{rem: corrections} (new): corrections to \cite{EKYY1} and \cite{EKYY2}, see ``Corrections to the cited works'' below.
\end{itemize}

\subsection*{Section 4}
\begin{itemize}
\item Beginning of Section 4, proof of the necessary condition, (4.1)--(4.8): ``$\lim_{s\to\infty}s^4\p(|x_{12}|\ge s)\neq0$'' $\to$ ``$s^4\p(|x_{12}|\ge s)$ does not tend to $0$''; $r_k\ge4$; in (4.2) $|h_{ii}|\le1$ (was $<1$); explicit $c_2=c_1/1024$, $c_3=e^{-c_2/4}$ and $N\ge\max(2,2\E|x_{11}|^2)$; the independence used in (4.8) is justified through the first index pair with $|h_{ij}|\ge4$.
\item \eqref{suff}: additive error $N^{-\delta}$ $\to$ $\epsilon_N\to0$, with the quantifiers and dependence of \eqref{def: ranvec}; ``this is \eqref{def: ranvec}''.
\item Remark \ref{rem: additive error} (new): for the law of Remark \ref{example}, \eqref{suff} fails with $\epsilon_N=N^{-\delta}$ for every $\delta>0$; this example is not formalized; the error $N^{-\delta'}$ under higher moments is expected and not proved here. The paragraph after it: the comparison matrix is $\wt H^S$, bounded almost surely.
\item Proof of the sufficient condition, \eqref{de61}--(4.13): ``for fixed $\epsilon>0$'' $\to$ ``fix $0<\epsilon<1/30$'', with the step that needs each bound; (4.12)--(4.13) in $o(\cdot)$ form (the printed $\delta$ clashed with the $\delta$ of \eqref{suff}) with a proof; the conditional means $m$, $\wt m$.
\item \eqref{95}--\eqref{jlxf}: laws defined by densities shifted by $\beta$ $\to$ conditional laws centred by $m$ (no densities); \eqref{de67} is an identity on a product space with the law of $H_N$ (was ``$\sim$''); $M=N^{-1/2}mJ+N^{-1/2}(\wt m-m)I$ and \eqref{jlxf} bounds the shift of $\lambda_N$.
\item (4.18)--(4.19): written as integrals against $\nu$ and $\wt\nu$; second moment $1-o(N^{-1+2\epsilon})$; diagonal second moment $\le C_1$.
\item \eqref{qbtt}: $H^S$ $\to$ $\wt H^S=\sigma^{-1}H^S$, bounded almost surely by $N^{-\phi_s}$, $\phi_s=3\epsilon/4$; Theorem \ref{thm: new edge un} applied in the orientation of \eqref{qset1}.
\item (4.21)--(4.25): ``$>3/4$'' with a proof; in (4.22) the indicator of $|\wt H_{ij}H^C_{ij}|\le3/4$ (printed without absolute value); the event SelGood and $\wt H^C$ defined explicitly; \eqref{mkgd} proved; $H=H^S+E+M$ on an event of probability $1-o(1)$.
\item \eqref{segg}--(4.37): \eqref{segg} is stated for $\wt H^S$ and $\wt E=\sigma^{-1}E$; the derivation of \eqref{suff} from \eqref{segg} is written out on the scale $N^{2/3}(\lambda_N-2)$, with $\delta=\frac12\min(\delta_0,1/24)$ and $\epsilon_N=N^{-\delta_0}+p_N$.
\item (4.27)--\eqref{FG3}: only the nonzero positions of $\wt E$ are kept, so $D$ is invertible and $V$ has orthonormal columns; \cite[Lemma 6.1]{KY} is replaced by the matrix determinant lemma, with its proof (\eqref{tyyz} and the identity after it).
\item \eqref{928o}: maximum over $1\le i\le N$ (was $1\le i\le 2s+t$), bound $\varphi^C(N^{-\phi_s}+N^{-1/3})\le N^{-\epsilon/2}$.
\item \eqref{928}: maximum over the index set touched by $\wt H^C$, a function of $H^C$ only; proved from Lemma \ref{lem: WBGij} by Markov's inequality.
\item (4.40)--(4.47): $\wt\mu:=\lambda_N\pm N^{-3/4}$ $\to$ the one-sided barrier $\lambda_N+N^{-3/4}$; (4.41) ``the largest eigenvalues of GOE are separated on the scale $N^{-2/3}$'' (with Remark \ref{genedge}) $\to$ $\lambda_\alpha\le\wt\mu-N^{-3/4}$, so \eqref{ttg} holds deterministically; \eqref{ttg2} from \eqref{rigiditynew} (was ``see \eqref{qset1}''); (4.44) ``$=$'' $\to$ ``$\le$''; \eqref{2ss} $CN^{-1/4}$ $\to$ $\varphi^CN^{-1/4}$; (4.46) a Gershgorin bound with $\epsilon<1/30$; (4.47) the upper bound by continuity in $\gamma$ and the lower bound by the Rayleigh quotient, $\mu_N\in[\lambda_N-\varphi^CN^{-1+5\epsilon},\lambda_N+N^{-3/4})$.
\end{itemize}

\subsection*{Section 5}
\begin{itemize}
\item Beginning of Section 5: $\wt H$ ``satisfies the bounded support condition with $q=O(N^{1/2}/\log N)$ and has the same first four moments'' $\to$ four moments off the diagonal, two on it, entries bounded almost surely by $\wt q^{-1}$ with $\wt q=N^{1/2}/((2C_1+4)C_1\log N)$.
\item Lemma \ref{lem: match}, (5.1), and its proof: hypotheses are independence and the moment conditions of Theorem \ref{thm: new rigidity} (no support condition); explicit $\wt q$; $\wt H$ satisfies \eqref{eq:diag} with $\max(C_1,1)$; ``the diagonal part is trivial'' $\to$ an explicit construction.
\item Lemma \ref{lem: four moments}, \eqref{tyss}--\eqref{tyss2}, and its proof (after Remark \ref{rem: four moments}), (5.4)--(5.9): ``$C\le|A|$'' $\to$ ``$|A|\le C$''; $\E x^3=C$ $\to$ $\E x^3=A$; $D=2C+4$; ``$x=X+Y_t$, $t=4B-8A^2-2$'' (variance $2$) $\to$ $x=(X+Y_t)/\sqrt2$, $t=4B-8A^2-7$, for $B\ge2A^2+2$, and a mixture with the two-point law $\mu_A$ for $A^2+1\le B<2A^2+2$ (was ``it is easy to find'').
\item Paragraph after (5.9), before Lemma \ref{lem:p moment}: new term ``matched pair''.
\item Lemma \ref{lem:p moment}, \eqref{asXY}--\eqref{sqbm}: explicit quantifiers ($C_*$ depends on $(C_1,\phi')$; for all $c',C'$ there are $C$, $N_0$); matched pairs with $q\ge N^{\phi'}$; \eqref{asXY} is assumed to fail with probability at most $N^{C'}e^{-c'\varphi^3}$; \eqref{sqbm} for every even $p\le\varphi$.
\item Lemma \ref{lem:green}, \eqref{jhw2}--\eqref{sbc22}: for every $C_F$ there are $\wt\epsilon$, $C$; for every $\epsilon<\wt\epsilon$ there is $N_0$; $F\in C^3$ with \eqref{jhw2}; \eqref{mainpear} for all $E\in I_\epsilon$ and \eqref{sbc22} for all $E_1,E_2\in I_\epsilon$ (oriented).
\item Proof of Theorem \ref{thm: new rigidity}, \eqref{tby}--\eqref{tby3}: level $\zeta=4$ with $a=49+3C_s+C_*$ (printed $a\ge12$, which is not enough at this level); \eqref{3737T} with the window $E\le3$ (printed $E\le2+q^{-1}$) and offset $K$; $p=2\lfloor\varphi/2\rfloor$ (printed $p=\varphi$); the criterion $\im m<(2N\eta)^{-1}$ for no eigenvalue near $E$ (printed ``$\ll$''); (5.22) $\varphi^{C_t/2-K-a}N^{-1/3}\le(N\eta)^{-1}\le\varphi^{-K-a}$ (printed $\varphi^{O(1)}N^{1/3}\le(N\eta)^{-1}\le\varphi^{O(1)}N^{-\phi/2}$); a grid in $E$; \eqref{tby3} by reflection.
\item \eqref{tby4} and the last step of the proof of Theorem \ref{thm: new rigidity}: \eqref{tby4} from Lemma \ref{lem: old} at $\zeta=3$ for $\wt H$ and a grid; the last step follows \cite[(5.4)--(5.15)]{EYY} with explicit parameters and uses the norm bound \eqref{tby3}.
\item Lemma \ref{lem: edge un str}, \eqref{hh1}--\eqref{qset15}, and its proof: ``two generalized Wigner matrices'' satisfying \eqref{rigiditynew}, \eqref{rigiditynew2}, \eqref{tby4}, \eqref{hh1}, \eqref{hh2} and \eqref{jhw2}--\eqref{sbc22} $\to$ pairs of matrices with independent entries satisfying (1) $\|H\|\le3$, \eqref{tby4} for $\eta\le N^{-2/3}$, \eqref{hh1} and \eqref{hh2} with high probability for every $C\ge C_0$, and (2) a comparison hypothesis with test functions with four derivatives and rate $N^{-\kappa+C'\epsilon}$; \eqref{qset15} for all $s\in\mathbb R$ (was $s>0$); the proof is written in six steps, with the corrections to \cite{EYY} (the sentence that lists \cite[(2.19), (2.25), (2.26), (6.2), (6.3)]{EYY} as the hypotheses is removed).
\item Proof of Theorem \ref{thm: new edge un} after Lemma \ref{lem: edge un str}, (5.28): the hypotheses of Lemma \ref{lem: edge un str} are checked for $H^{\V}$ and for $\wt H$, and Lemma \ref{lem: ex edge un} is applied to $\wt H$ at $\phi=2/5$.
\end{itemize}

\subsection*{Section 6}
\begin{itemize}
\item Beginning of Section 6, \eqref{813}--\eqref{813tr}: the class of pairs of the section is stated (independent entries, matched moments, $|h_{ij}|\le q^{-1}$ almost surely, $q\ge N^\phi$; for Theorem \ref{thm: new rigidity}, $q=\min(N^\phi,\wt q)$ and $\phi$ replaced by $\min(\phi,1/4)$); ``bounded support condition with $q=N^{-\phi}$'' removed; all $H_\gamma$ on one probability space; \eqref{813} for each $\gamma$ with constants uniform in $\gamma$, proved from Lemma \ref{lem: old} (the argument ``finitely many distributions'' is removed); ${\bf k}\in\mathbb R^n$ $\to$ $\mathbb N^n$.
\item Example 6.1, discussion after (6.5) and (6.7): ``bounded by $N^{-\phi}$ with high probability'' $\to$ ``almost surely''; ``the first four moments of $H$ and $\wt H$ coincide'' $\to$ ``the first four moments of the off-diagonal entries coincide (in this example $(a,b)$ is off-diagonal)''.
\item Example 6.1, the sum in the heuristic after (6.5): $\sum_{k=0}^4A_k$ $\to$ $\sum_{k=0}^mA_k$ (misprint).
\item Definition \ref{taz}, \eqref{taz1}--\eqref{PPP}, and Lemma 6.4, \eqref{bG4}: $k\in\mathbb R$, ${\bf k}\in\mathbb R^s$ $\to$ $\mathbb N$, $\mathbb N^s$; ``Definition \ref{taz2}'' $\to$ ``Definition \ref{taz}''.
\item Lemma \ref{lem:S_expansion}, \eqref{eq:S_exp2}--\eqref{boundcalA2}, and the sentence before it: $H_{\gamma'}$ any of the interpolating matrices and $(a,b)$ arbitrary; ``$s=O(\varphi)$, $\zeta=O(\varphi)$'' $\to$ $1\le\zeta\le\varphi^2$, $s\le K_s\varphi^2$, $\xi=\lceil2\zeta/\phi\rceil$, with $N_0$ after $(C_1,\phi,K_s)$; the matched order $m_0=4$ off the diagonal and $m_0=2$ on it (printed $4$ and $5$); real coefficients independent of $z$ and of the indices; \eqref{boundcalA} in two cases ($N^{-2-|\alpha|\phi/10}$ and $N^{-1-|\alpha|\phi/4}$); \eqref{eq:S_exp3} with $\wt{\mathcal P}$ and the same coefficients; the estimates hold for products with conjugate and centred factors; remainder $4^sN^{-\zeta}$.
\item The iteration, (6.26)--\eqref{bssm}: rewritten with words, formal sums and data $d$; $1\le\zeta\le\varphi$, $s\le K_s\zeta$; diagonal swaps match two moments; per-step gain $N^{-\phi/5}$, absorption \eqref{absorb} (new label) and induction \eqref{722s}; the termination ``$\sum|{\bf k}_i|\ge5n$'' $\to$ $16\cdot4^sN^{-6\zeta/5}\le N^{-\zeta}$; \eqref{HHtt} is the sum form with explicit weights and \eqref{conK} the admissible data; the max forms \eqref{HHtt2}, \eqref{HHtt3} with the coefficient function \eqref{fIJ} (new label), \eqref{bssm} for centred words through \eqref{eq:S_exp3}.
\item Lemma \ref{lem:words} (new): a bound on words with explicit constants.
\item Proof of Lemma \ref{lem: WBGij}, \eqref{tddc}--\eqref{sysz}: the case $\eta>10$ by the Schur complement; $\wt q\sim N^{1/2}/\log N$; the proof uses the sum form \eqref{HHtt} with $\min(\phi,1/4)$ and explicit weights (the max form does not give \eqref{res: WBGij}); \eqref{tddc} keeps $\wt q^{-1}$.
\item Proof of Lemma \ref{lem:p moment}, (6.42)--(6.45): \eqref{sqbm} is proved directly for words with conjugate factors; the max form with $s=\zeta=p$; block structure \eqref{dsb}; block bound $3^{|\beta|}Q$ with $Q=X^2+Y+18N^{-1}$.
\item Proof of Lemma \ref{lem:S_expansion}, first paragraph (after (6.45)) and \eqref{defxin}: the matched order $m_0$; the entries of $H_{\gamma'}$ and $Q$ are bounded by $q^{-1}$ almost surely; $\xi:=\lceil2\zeta/\phi\rceil$.
\item Proof of Lemma \ref{lem:S_expansion}, (6.47)--\eqref{hxxt}: the diagonal case of $V$; new paragraph on products with conjugate and centred factors (each factor expands separately, and the coefficients do not depend on the kinds of the factors); ``\eqref{eq:S_exp3} can be proved analogously'' $\to$ the proof with $\wt{\mathcal P}$; $I^a_{b,c}\subset\mathbb N^a$.
\item Proof of Lemma \ref{lem:S_expansion}, (6.57)--(6.72): $4$, $5$ $\to$ $m_0$, $m_0+1$ in the sums; the bad event and the independence of $R$ and $h_{ab}$ are written out.
\item \eqref{wbtnm}: the third summation variable ${\bf k_2}$ $\to$ ${\bf k_3}$; $\E$ inserted before $\mathcal P$; $k_{m+1,t}$ $\to$ $k_{m+1}(t)$ (misprints).
\item Proof of Lemma \ref{lem:S_expansion}, \eqref{off3}--(6.79): from \eqref{off3} to \eqref{eq:S_exp2} with an explicit formula for $\mathcal A_\alpha$; (6.78) with the constant $C_1$ restored and a diagonal case; \eqref{boundcalA} in both cases; \eqref{eq:S_exp4} and \eqref{boundcalA2} written out.
\item Proof of Lemma \ref{lem:green} (after (6.79)), (6.80)--(6.105): Lemma \ref{lem:S_expansion} and \eqref{bssm} are used at $\phi'=\min(\phi,1/4)$; (6.81) with the constant $2$; (6.83) $N^{-\min(\phi,1/3)+C\epsilon}$; in \eqref{eq:telescopic_fn} the second term $\wt G_{ij}\wt G_{ji}$ $\to$ $\wt G_{ij}\overline{\wt G_{ij}}$; \eqref{eq:taylor} with an explicit remainder; the diagonal swaps add $N\cdot N^{-1-\phi+C\epsilon}$; \eqref{yrsh} $\xi=\lceil3/\phi\rceil$; in \eqref{pqbh0} $C$ does not depend on $p$; \eqref{pqbh}--\eqref{pqbh3} with window $\lceil4p/\phi'\rceil$; in \eqref{Gabcd} the third form $G_{i_ta}G_{bj_t}$ $\to$ $\wt G_{i_ta}\wt G_{bj_t}$; the extra ``(5.10)'' sentence removed; \eqref{yrsh2} $R_{ij}R_{ji}$ $\to$ $R_{ij}\overline{R_{ij}}$; in (6.105) $F^{(s)}(x_R)$ inside $\E$ and ${\bf k}_t\in I^2_{3/\phi,|{\bf k}_t|}$.
\end{itemize}

\subsection*{Misprints}
The misprints corrected above are: $\bigcup\to\bigcap$ in \eqref{scm}, \eqref{Gij estimate}, (3.6) and \eqref{rigiditynew}; the direction of the bound on $\eta$ before \eqref{tyzz}; $\mathbb R\to\mathbb N$ in Section 6; the absolute value in (4.22); the window $E\le3$ in \eqref{3737T}; the conjugates in \eqref{eq:telescopic_fn} and \eqref{yrsh2} and the position of $\E$ in (6.105); the tildes in \eqref{Gabcd}; the indices and the $\E$ in \eqref{wbtnm}; and the upper limit of the sum in the heuristic of Example 6.1. The level-$\zeta$ form of \eqref{tyzz}, the norm bound used in the last step of the proof of Theorem \ref{thm: new rigidity}, the exponent $\min(\phi,1/4)$ in the proof of Lemma \ref{lem:green}, and the products with conjugate and centred factors in Lemma \ref{lem:S_expansion} are corrections of statements or arguments, listed above.

\subsection*{Corrections to the cited works}
The corrections to \cite{EKYY1} and \cite{EKYY2} found in the formalization are listed in Remark \ref{rem: corrections}. The corrections to \cite{EYY} are given in the proof of Lemma \ref{lem: edge un str}. None of them changes a statement used in this paper.
\endgroup
